\section{引言与深度学习基础回顾}

本节课是 Stanford CS224N 的第五讲，由 Chris Manning 教授主讲。课程内容分为三个主要部分：（1）深度学习实用技巧的补充（正则化、参数初始化、优化器等）；（2）语言模型（Language Models）的概念与构建方法；（3）循环神经网络（Recurrent Neural Networks, RNN）的原理与应用。

\begin{figure}[H]
\centering
\includegraphics[width=0.95\textwidth]{slides-images/slide-000.jpg}
\caption{CS224N Lecture 5 封面：Language Models and Recurrent Neural Networks\protect\footnotemark}
\end{figure}
\footnotetext{Slides 第1页。}

\subsection{从浅层到深层：深度学习的崛起}

Manning 教授首先回顾了神经网络的发展历史。在 1980--1990 年代，虽然反向传播算法已经被 Hinton 等人推广，但当时几乎所有的神经网络都只有\textbf{一个隐藏层}。原因在于，人们长期无法让深层网络比浅层网络表现更好。

\begin{figure}[H]
\centering
\includegraphics[width=0.95\textwidth]{slides-images/slide-005.jpg}
\caption{Bengio et al. (2006) 指出：训练深层神经网络曾被认为过于困难，深层网络的表现通常不如一两个隐藏层的浅层网络\protect\footnotemark}
\end{figure}
\footnotetext{Slides 第6页。}

\begin{knowledgebox}{深度学习的历史转折}
直到 2006 年前后，Yoshua Bengio 等人提出逐层贪心预训练（Greedy Layer-Wise Training）等方法，深度网络才开始逐步展现出优于浅层网络的能力。Manning 指出，推动这一转变的并非某个惊天发现，而是若干``小想法''和``小技巧''的积累——正则化方法、激活函数选择、参数初始化策略、优化器改进等。这些看似不起眼的改进，让一个停滞了 15 年的领域焕发新生。
\end{knowledgebox}

\subsection{正则化：从经典到现代视角}

\subsubsection{经典正则化}

正则化（Regularization）是通过在损失函数中添加额外约束来防止过拟合的技术。最常见的形式是 L2 正则化：

$$J(\theta) = J_{\text{data}}(\theta) + \lambda \sum_i \theta_i^2$$

\begin{itemize}
    \item $J_{\text{data}}(\theta)$：数据拟合项（如交叉熵损失）
    \item $\lambda \sum_i \theta_i^2$：正则化项，鼓励参数保持较小的值
    \item $\lambda$：正则化强度超参数
\end{itemize}

\begin{figure}[H]
\centering
\includegraphics[width=0.95\textwidth]{slides-images/slide-003.jpg}
\caption{经典的过拟合视角：训练误差持续下降，但验证误差在某一点后开始上升\protect\footnotemark}
\end{figure}
\footnotetext{Slides 第4页。}

\subsubsection{现代视角：Double Descent}

Manning 特别强调，现代大规模神经网络的行为与经典理论\textbf{截然不同}。在足够大的网络中，即使模型完全记住了训练数据（训练损失趋近于零），验证损失仍然会继续下降。

\begin{importantbox}{现代正则化的范式转变}
现代大型神经网络可以将训练数据几乎完美拟合（训练损失接近零），这在经典统计学看来是灾难性的过拟合。但实际上，只要正则化做得好，模型仍然能很好地泛化到新数据。这一发现被称为``Double Descent''现象，彻底改变了我们对模型复杂度和泛化的理解。
\end{importantbox}

\begin{figure}[H]
\centering
\includegraphics[width=0.95\textwidth]{slides-images/slide-004.jpg}
\caption{现代视角：大型网络中，验证损失先升后降，最终随着训练继续而改善\protect\footnotemark}
\end{figure}
\footnotetext{Slides 第5页。}

\subsubsection{Dropout：现代神经网络的标配正则化}

传统的 L2/L1 正则化对大型神经网络往往不够强。Manning 介绍了 \textbf{Dropout}——一种专为神经网络设计的正则化方法。

\begin{figure}[H]
\centering
\includegraphics[width=0.95\textwidth]{slides-images/slide-006.jpg}
\caption{Dropout 示意图：训练时随机将部分神经元置零，测试时使用全部神经元但进行缩放\protect\footnotemark}
\end{figure}
\footnotetext{Slides 第7页。}

\begin{importantbox}{Dropout 的核心机制}
\begin{enumerate}
    \item \textbf{训练时}：对每个训练样本，随机生成一个由 0 和 1 组成的掩码（mask），通过 Hadamard 积将部分神经元的输出置零。每次使用不同的掩码。
    \item \textbf{测试时}：不进行 Dropout，保留所有权重，但对输出进行缩放以补偿训练时的丢弃。
\end{enumerate}
\end{importantbox}

Dropout 为什么有效？Manning 给出了三种直觉解释：

\begin{itemize}
    \item \textbf{防止特征共适应}（Feature Co-adaptation）：模型不能过度依赖某几个特征的组合，因为它们随时可能被删除
    \item \textbf{隐式模型集成}：每次 Dropout 相当于训练一个不同的子网络，最终效果类似于指数级数量的模型集成
    \item \textbf{介于 Naive Bayes 和 Logistic Regression 之间}：Naive Bayes 独立地评估每个特征，Logistic Regression 在所有特征的上下文中设置权重，Dropout 则在随机子集的上下文中学习，提供了一种中间地带
\end{itemize}

\subsection{向量化：深度学习的性能基石}

\begin{warningbox}{永远不要写 for 循环}
深度学习的成功很大程度上依赖于向量化计算。在 Python 中，使用 for 循环逐元素处理数据比向量/矩阵操作慢一到两个数量级。在 GPU 上，这个差距更是达到两到三个数量级。如果你发现自己在写 for 循环处理数据（而不是最表层的输入预处理），几乎一定是做错了。
\end{warningbox}

\subsection{参数初始化}

神经网络的参数初始化至关重要。如果将所有参数初始化为零或同一常数，网络会陷入\textbf{对称性陷阱}——所有神经元学到的东西完全相同，相当于只有一个神经元。

\begin{figure}[H]
\centering
\includegraphics[width=0.95\textwidth]{slides-images/slide-010.jpg}
\caption{参数初始化：必须使用小的随机数初始化权重，避免对称性问题。Xavier 初始化根据层的输入输出维度自适应设置方差\protect\footnotemark}
\end{figure}
\footnotetext{Slides 第11页。}

\begin{importantbox}{Xavier 初始化}
Xavier/Glorot 初始化设置权重的方差为：
$$\text{Var}(W_i) = \frac{2}{n_{\text{in}} + n_{\text{out}}}$$
其中 $n_{\text{in}}$ 是前一层的大小，$n_{\text{out}}$ 是后一层的大小。这确保信号在前向和反向传播过程中既不会消失也不会爆炸。Manning 指出，后来 Layer Normalization 的引入在很大程度上消除了对精确初始化的需求。
\end{importantbox}

\subsection{优化器：从 SGD 到 Adam}

Manning 简要介绍了优化器的发展：

\begin{figure}[H]
\centering
\includegraphics[width=0.95\textwidth]{slides-images/slide-011.jpg}
\caption{各种优化器及其关系\protect\footnotemark}
\end{figure}
\footnotetext{Slides 第12页。}

\begin{itemize}
    \item \textbf{SGD}（随机梯度下降）：最基础的优化器，理论上能解决几乎所有问题，但对学习率和调度非常敏感
    \item \textbf{Adagrad}：为每个参数维护历史梯度平方和，自适应调整学习率。缺点是学习率会过早衰减
    \item \textbf{Adam}：结合动量和自适应学习率，是最常用的默认选择
\end{itemize}

\begin{knowledgebox}{如果只记住一个优化器}
Adam 是最安全的默认选择。它对超参数不敏感，几乎在所有场景下都能给出不错的结果。对于词向量等稀疏更新场景，带有权重衰减的变体（AdamW）可能更合适。
\end{knowledgebox}

\subsection{本章小结}

\begin{itemize}
    \item 深度学习的成功建立在若干关键技巧的积累之上：Dropout 正则化、合理的参数初始化、自适应优化器
    \item 现代大型网络的过拟合行为与经典理论不同——Double Descent 现象表明，足够大的模型即使完美拟合训练数据仍能良好泛化
    \item 向量化计算是深度学习效率的基石，应避免在数据处理中使用 for 循环
    \item Adam 优化器是实践中最可靠的默认选择
\end{itemize}

%% ========== Part 2: 语言模型 ==========

